\chapter{Redefinición del Alcance y Objetivos del Proyecto} \label{sec:alcance}
Durante el transcurso del proyecto se identificaron una serie de restricciones de carácter económico, logístico y operativo que, en conjunto, tornaron inviable la ejecución del plan original en su totalidad. A continuación se describen los factores determinantes que motivaron la redefinición del alcance.

El plan original contempló el desarrollo de tres diseños de circuito impreso distintos: el amplificador de operación permanente en Clase AB, el amplificador de cresta en Clase C y los elementos pasivos de la topología Doherty. El costo unitario del transistor empleado, sumado a los costos de fabricación, envío y ensamblado de cada iteración de placa, representó una restricción económica significativa. Si bien las demoras asociadas al proceso de fabricación y envío desde el exterior fueron contempladas en la planificación original, la acumulación de estos tiempos a lo largo de tres iteraciones de fabricación habría comprometido seriamente el cronograma del proyecto. Adicionalmente, el ensamblado del primer circuito impreso fue realizado sin costo por parte de la empresa Outsource, beneficio que no podía garantizarse para las iteraciones subsiguientes, lo que introdujo una incertidumbre adicional tanto en los costos como en los plazos.

Desde el punto de vista del acceso a equipamiento de medición, el departamento de ingeniería electrónica de la facultad presentó limitaciones tanto de alcance como de funcionamiento de su instrumental disponible para la caracterización de circuitos de radiofrecuencia a las frecuencias de trabajo consideradas en este proyecto. Por un lado, el analizador de espectro de mayor rango disponible en la facultad alcanza una frecuencia máxima de 3~GHz, insuficiente para la medición del contenido armónico del amplificador, dado que el tercer armónico de la frecuencia de operación se ubica en 6,6~GHz, fuera del rango de dicho instrumento. Por otro lado, se detectó durante las primeras sesiones de medición que el tono generado a la salida del generador de señales de RF disponible no coincidía con la frecuencia indicada en el equipo, lo que comprometió adicionalmente la posibilidad de realizar mediciones confiables con el equipamiento propio de la institución. Para subsanar estas limitaciones se recurrió a la Comisión Nacional de Energía Atómica (CNEA) y a la Universidad Nacional de San Martín (UNSAM), instituciones donde finalmente se llevaron a cabo las mediciones. Sin embargo, el acceso a los equipos y a la ayuda por parte del personal en dichas instituciones requirió coordinación anticipada y estuvo sujeto a la disponibilidad tanto de los integrantes del grupo como de las propias instituciones, lo que limitó no solo la frecuencia con que pudieron realizarse sesiones de medición, sino también el tiempo disponible dentro de cada sesión, dado que los equipos y los espacios de trabajo debieron compartirse con otras actividades propias de cada institución. Extender este esquema a la caracterización de tres diseños distintos habría requerido una cantidad y duración de sesiones de medición difícilmente compatible con los plazos disponibles.

En relación con los ensayos de irradiación contemplados en el anteproyecto, se presentaron dificultades para dar con una instalación experimental disponible cuyo tamaño de muestra admisible fuera compatible con las dimensiones del primer circuito impreso fabricado. Esta incompatibilidad física, sumada a la complejidad logística asociada a la coordinación de turnos de irradiación, llevó a desestimar la realización de este ensayo dentro del alcance del presente trabajo.

Frente a este conjunto de restricciones, y considerando que los resultados obtenidos en la caracterización del amplificador Clase AB, satisficieron en gran medida los objetivos de prestaciones establecidos en el anteproyecto, se tomó la decisión de redefinir el alcance del proyecto. El desarrollo quedó acotado al amplificador de potencia en Clase AB y al acoplador híbrido en cuadratura. El amplificador Clase AB constituyó por sí mismo un amplificador de potencia funcional que cumplió con gran parte de los requisitos de la misión, mientras que el acoplador híbrido representó una primera instancia de desarrollo que sentó las bases para una eventual continuación del trabajo hacia la implementación completa de la topología Doherty en el futuro. El ensayo de irradiación fue excluido del alcance actual, redirigiendo los esfuerzos hacia una caracterización experimental más exhaustiva de los parámetros del amplificador fabricado.